\section{\ssdos and the machine}\label{sec:ssdos}

% \ssdos is a real-time system for solid state drives using a tickbased pipeline
% for program execution targeting risc-v.

\ssdos is an operating system for solid state drives running on risc-v
processors.
Programs running on the system is divided into ticks in a pipeline, where each
tick is a fixed time interval.

One of the target machines, MicroBlaze-V\cite{sigma} use --- instead of caches
for its memory hierarchy --- scratch and shared memory.
For a multi processor system, this means that each processor has its own
scratch memory. Being wired specifically for each processor, this gives faster
memory speeds. The other memory type is shared, of which each processor is
wired to. This is however slower.
% Each memory region is accessed by a specific memory address, as expected.

Allocation on this system is done via application specific arena allocators,
such as in figure \ref{code:arena} with an accompanying allocation function in
figure \ref{code:alloc}. A simpler allocator simply returns the $\texttt{start}
+ \texttt{allocated}$  pointer and increases the \texttt{allocated} field by
how much is being requested.
This is essentially a stack.


% We are not really considering a heap at all (even though obviously program text
% exists, and \var{.words} can exist there, too).

% \subsection{Model}
% This project specifically targets the cacheless \texttt{v80}
% platform\cite{sigma}. It consists of four clusters with four AMD MicroBlaze V
% processors, each having its own scratchpad memory, and all clusters connected
% to a shared scratchpad, referred to as shared moving forwards.
% The key take-away is that shared memory is slower than scratch memory.
% 
% \subsection{Allocating}
% 
\begin{figure}
    \begin{lstlisting}[language=C, gobble=4]
        struct memory_region
        {
            uint8_t *start;
            size_t   size;
        };

        struct arena_allocator
        {
            struct memory_region region;
            size_t               allocated;
        };
    \end{lstlisting}
    \caption{Structure of an allocator}
    \label{code:arena}
\end{figure}

\begin{figure}
    \begin{lstlisting}[language=C, gobble=8]
        uint8_t *arena_alloc(struct arena_allocator *alloc, size_t size, size_t alignment)
        {
            uint8_t     *free_region_start = alloc->region.start + alloc->allocated;
            const size_t alignment_padding =
                    (alignment - (((size_t) free_region_start) % alignment)) %
                    alignment;
            const size_t size_to_alloc = size + alignment_padding;

            const size_t remaining =
                    (alloc->region.start + alloc->region.size) - free_region_start;

            if (remaining < size_to_alloc)
                return NULL;

            uint8_t *allocated = free_region_start + alignment_padding;
            alloc->allocated += size_to_alloc;
            return allocated;
        }
    \end{lstlisting}
    \caption{Example of an allocator function}
    \label{code:alloc}
\end{figure}
% 
